\chapter{Formy kwadratowe i przeszkoda chirurgiczna}
\label{ch:przeszkoda-chirurgiczna}
\index{grupa L@Grupa $L$}

W Rozdziale~\ref{ch:chirurgia-geometryczna-21} operacja na
obramowanej sferze zmieniała rozmaitość, a forma przecięć
na jądrze mapy stopnia jeden pokazywała, dlaczego w
\emph{środkowym wymiarze} nie można usuwać klas po jednej.
Teraz zapiszemy to ograniczenie algebraicznie. Otrzymamy
grupę, w której klasa zerowa oznacza możliwość dokończenia
chirurgii przy odpowiednich założeniach wymiarowych.
Odróżnimy twierdzenia o konkretnych macierzach, które
udowodnimy tutaj, od głębokiego twierdzenia chirurgicznego
Walla, którego geometryczne lematy przywołamy jawnie.

\section{Drogi przecięcia i pierścień grupowy}
\label{sec:pi-formy-22}

Jeśli $X$ jest spójna, połóżmy $\pi=\pi_1(X)$ i
$\Lambda=\Z[\pi]$. Element $\Lambda$ jest skończoną
sumą $\sum_g n_g g$, gdzie $n_g\in\Z$. Mnożenie
przedłuża liniowo mnożenie w $\pi$. W przypadku
$\pi=1$ otrzymujemy zwykły pierścień $\Z$.
Niech $w:\pi\to\{\pm1\}$ będzie charakterem
orientacji celu. Na $\Lambda$ wprowadzamy \emph{inwolucję}
\begin{equation}\label{eq:inwolucja-L-22}
 \overline{\sum_g n_g g}
     =\sum_g n_g w(g)g^{-1}.
\end{equation}
Odwrócenie drogi zamienia $g$ na $g^{-1}$, a
$w(g)$ mówi, czy przeniesienie lokalnej orientacji
wzdłuż tej drogi zmienia znak.

\begin{stwierdzenie}[Dlaczego to inwolucja]
\label{prop:inwolucja-L-22}
Przepis~\eqref{eq:inwolucja-L-22} spełnia
$\overline{ab}=\bar b\bar a$ i
$\overline{\bar a}=a$.
\end{stwierdzenie}
\begin{proof}
Wystarczy sprawdzić elementy grupy, ponieważ obie
strony są addytywne. Ponieważ $w$ jest homomorfizmem,
\[
 \overline{gh}=w(g)w(h)h^{-1}g^{-1}
              =(w(h)h^{-1})(w(g)g^{-1})
              =\bar h\bar g.
\]
Ponadto $w(g^{-1})=w(g)$ i $w(g)^2=1$, skąd
$\overline{\bar g}=g$. Przedłużenie liniowe kończy
rachunek.
\end{proof}

Po podniesieniu sfer do nakrycia uniwersalnego
pojedynczy punkt ich przecięcia wnosi do
\emph{współczynnika przecięcia} znak razy element $g\in\pi$.
Znak zależy od orientacji, a $g$ od obranych dróg
od punktu bazowego do sfer. Dlatego dwa przecięcia
o przeciwnych znakach mogą sumować się do
$g-h\ne0$ w $\Lambda$, choć ich zwykła liczba
przecięcia jest zerem. Zmiana dróg zmienia zapis
macierzy, lecz nie jej klasę izometrii.

\begin{definicja}[Forma $\epsilon$-hermitowska]
\label{def:hermitowska-22}
Niech $P$ będzie skończenie generowanym projektowym
lewym modułem
nad $\Lambda$ oraz $\epsilon\in\{\pm1\}$.
\emph{Formą $\epsilon$-hermitowską} nazywamy
odwzorowanie $\lambda:P\times P\to\Lambda$ takie, że
\[
 \lambda(ax,by)=a\lambda(x,y)\bar b,\qquad
 \lambda(y,x)=\epsilon\overline{\lambda(x,y)}.
\]
Jest \emph{nieosobliwa}, gdy jej odwzorowanie
sprzężone $P\to P^*$ jest izomorfizmem. Dual
ma tu strukturę modułu z uwzględnieniem
inwolucji~\eqref{eq:inwolucja-L-22}.
\end{definicja}

Dla $n=2k$ przecięcia $k$-wymiarowych sfer
spełniają ten wzór z $\epsilon=(-1)^k$:
zamiana kolejności dwóch przestrzeni stycznych
w punkcie przecięcia zamienia bloki po $k$ wektorów,
co daje znak $(-1)^{k^2}=(-1)^k$; odwraca także
drogę, stąd kreska nad współczynnikiem.

\section{Samoprzecięcia i formy kwadratowe}
\label{sec:kwadratowa-22}

Sama $\lambda$ nie zapamiętuje, czy reprezentant
klasy daje się przedstawić jako \emph{obramowana,
osadzona} sfera. Potrzebna jest poprawka
samoprzecięciowa $\mu$; jej precyzyjną postać
podajemy w języku algebraicznym. Dla
$\epsilon=(-1)^k$ oznaczmy
\[
 Q_\epsilon(\Lambda)=
 \Lambda/\{a-\epsilon\bar a:a\in\Lambda\}.
\]
W mianowniku chodzi o podgrupę addytywną, a nie
zwykle o ideał. Nawias $[a]$ oznacza klasę $a$.

\begin{definicja}[Forma $\epsilon$-kwadratowa]
\label{def:forma-kwadratowa-22}
To para $(\lambda,\mu)$ na $P$, gdzie $\lambda$
jest nieosobliwą formą $\epsilon$-hermitowską, a
$\mu:P\to Q_\epsilon(\Lambda)$ spełnia
\begin{align}
 \mu(x+y)&=\mu(x)+\mu(y)+[\lambda(x,y)],
                   \label{eq:quad-add-22}\\
 \mu(ax)&=[a\mu(x)\bar a],                   \label{eq:quad-scal-22}\\
 \lambda(x,x)&=a+\epsilon\bar a
       \quad\text{dla dowolnego }a\in\mu(x).
                   \label{eq:quad-diag-22}
\end{align}
Ostatni wzór jest niezależny od reprezentanta
$a$: dodanie $c-\epsilon\bar c$ nie zmienia
$a+\epsilon\bar a$.
\end{definicja}

Geometrycznie $\mu(x)$ liczy samoprzecięcia
zanurzonej sfery i uwzględnia obramowanie.
Równanie~\eqref{eq:quad-add-22} mówi:
gdy połączymy reprezentantów $x$ i $y$, nowe
samoprzecięcia obejmują też przecięcia obu sfer.
W \eqref{eq:quad-diag-22} widać, dlaczego
„połowa przekątnej” nie zawsze wystarcza:
nad $\Z[\pi]$ nie można po prostu dzielić przez $2$.

\begin{stwierdzenie}[Rachunek przecięć na nakryciu]
\label{prop:rachunek-przeciec-Wall-22}
Niech $k\geq3$, a $M^{2k}\to X^{2k}$ będzie
$k$-spójnym odwzorowaniem normalnym stopnia jeden.
Wybierzmy obramowane, poprzeczne immersje
$k$-sfer reprezentujące $x,y\in K_k(M)$,
wraz z drogami od punktu bazowego. Każdy punkt
$p$ przecięcia wnosi znak $\varepsilon_p$ i element
$g_p\in\pi_1(X)$, wyznaczony przez przejście
między wybranymi podniesieniami sfer. Wtedy
\[
 \lambda(x,y)=\sum_{p\in x\pitchfork y}
                  \varepsilon_p g_p\in\Lambda.
\]
Wybierając przy każdym podwójnym punkcie immersji
$x$ jedną z dwóch kolejności gałęzi, otrzymujemy
\[
 \mu(x)=\left[\sum_{p\in\operatorname{Dbl}(x)}
                 \varepsilon_p g_p\right]
                 \in Q_{(-1)^k}(\Lambda).
\]
Wartości te nie zależą od regularnej homotopii
reprezentantów; spełniają
\eqref{eq:quad-add-22} i~\eqref{eq:quad-diag-22}.
\end{stwierdzenie}
\begin{proof}
Podnieśmy obie sfery do nakrycia uniwersalnego.
Współczynnik przy $g$ w $\lambda(x,y)$ jest
zwykłą algebraiczną liczbą przecięć jednego
podniesienia z przesunięciem drugiego przez $g$.
Jest więc stały przy homotopii. Po zamianie
obu sfer droga odwraca się, a znak zmienia się
o $(-1)^k w(g)$. To daje hermitowskość z
Definicji~\ref{def:hermitowska-22}.

Przy podwójnym punkcie zamiana kolejności
gałęzi zmienia $\varepsilon_p g_p$ na
$(-1)^k\overline{\varepsilon_p g_p}$.
Oba wyrazy mają tę samą klasę w
$Q_{(-1)^k}(\Lambda)$. Podczas ogólnej regularnej
homotopii punkty podwójne rodzą się lub znikają
parami o przeciwnych znakach i tej samej
etykiecie, zatem $\mu$ się nie zmienia.
Po połączeniu dwóch sfer rurką nowe punkty
podwójne to, oprócz dawnych, dokładnie ich
wzajemne przecięcia. Dostajemy stąd
\eqref{eq:quad-add-22}.

Na koniec odsuńmy immersję reprezentującą $x$
wzdłuż jej obramowania. Przy każdym podwójnym
punkcie odsunięta sfera przecina pierwotną
dwukrotnie, raz dla każdej kolejności gałęzi.
Ich suma wynosi $a+(-1)^k\bar a$ dla składnika
$a$ w $\mu(x)$. Wkład przekątnej z klasy Eulera
wiązki normalnej znika, bo sfera jest obramowana.
To jest~\eqref{eq:quad-diag-22}.
\end{proof}

Rachunek punktów wyjaśnia wzory, lecz aby
traktować $\mu$ jako funkcję na $K_k(M)$,
trzeba jeszcze wiedzieć, że dane normalne
wyznaczają właściwą klasę regularnej homotopii
obramowanych immersji. To identyfikacja
użyta w twierdzeniu 5.2 Walla; nie wynika
z samego zliczenia punktów.

Na przykład dwa przecięcia o znakach $+1$ i $-1$
oraz etykietach $1$ i $t$ dają
$\lambda(x,y)=1-t\in\Z[t,t^{-1}]$ dla
$\pi_1(X)=\langle t\rangle$. Zwykła liczba
przecięcia jest zerem, lecz okrąg złożony z
łuków obu sfer reprezentuje $t$ lub $t^{-1}$
i nie ogranicza
dysku Whitneya. Dopiero przeciwne znaki z
\emph{jednakową} etykietą pozwalają szukać
takiego dysku. Warunek $\mu(x)=0$ dopuszcza
analogiczne sparowanie samoprzecięć; przejście
od par do obramowanych dysków jest geometryczną
częścią twierdzenia 5.2 Walla.

\begin{przyklad}[Co zostaje, gdy $\pi=1$]
\label{prz:mu-Z-22}
Dla $\Lambda=\Z$ inwolucja jest zwykłą tożsamością.
Jeśli $k$ jest parzyste, $Q_+(\Z)=\Z$,
a~\eqref{eq:quad-diag-22} daje
$\lambda(x,x)=2\mu(x)$. Forma jest więc
\emph{parzysta}; $\mu(x)=\lambda(x,x)/2$.
Jeśli $k$ jest nieparzyste,
$Q_-(\Z)=\Z/2$; forma $\lambda$ jest
antysymetryczna i $\lambda(x,x)=0$,
ale $\mu(x)\in\Z/2$ może być niezerowe.
Ta ostatnia informacja nie daje się odczytać
z samej macierzy przecięć.
\end{przyklad}

\begin{definicja}[Para hiperboliczna i lagranżjan]
\index{para hiperboliczna i lagranzjan@Para hiperboliczna i lagranżjan}
\label{def:hiperboliczna-22}
Wolny moduł o bazie $(e_1,f_1,\dots,e_r,f_r)$
z macierzą w każdej parze
\[
 H_\epsilon=
 \begin{pmatrix}0&1\\\epsilon&0\end{pmatrix},
 \qquad \mu(e_i)=\mu(f_i)=0,
\]
oraz z zerowymi sprzężeniami między różnymi
parami, nazywamy \emph{hiperbolicznym}.
\emph{Lagranżjanem} formy kwadratowej jest
bezpośredni składnik $L\subset P$ taki, że
$L=L^\perp$ oraz $\mu|_L=0$.
W parze hiperbolicznej
$L=\langle e_1,\dots,e_r\rangle$ jest lagranżjanem.
\end{definicja}

Parę o macierzy $H_\epsilon$ realizują klasy
dwóch sfer czynnikowych w $S^k\times S^k$,
jak w Przykładzie~\ref{prz:forma-produkt-21}.
Lokalny model ich jednego przecięcia
pokazuje Rysunek~\ref{fig:para-hiperboliczna-22}.

\begin{figure}[htbp]
\centering
\begin{tikzpicture}[>=Latex,font=\small,line cap=round]
 \fill[manifoldblue!12,rounded corners=4pt]
   (-3.15,-.20) rectangle (3.15,.20);
 \fill[manifoldgold!19,rounded corners=4pt]
   (-.20,-1.12) rectangle (.20,1.12);
 \draw[manifoldblue!80!black,very thick]
   (-3.15,0)--(3.15,0);
 \draw[manifoldgold!85!black,very thick]
   (0,-1.12)--(0,1.12);
 \fill[manifoldred!85!black] (0,0) circle (2.5pt);
 \node[manifoldblue!80!black,anchor=west]
   at (-3.15,-.65) {$D^k\times\{0\}$};
 \node[manifoldgold!75!black,anchor=west]
   at (.38,.88) {$\{0\}\times D^k$};
 \node[manifoldred!85!black,anchor=west]
   at (.32,-.43) {jedno przecięcie};
\end{tikzpicture}
\caption{Lokalny model sfer
$e=S^k\times\{*\}$ i $f=\{*\}\times S^k$.
Ich dyski współrzędnych przecinają się
tylko w środku; po uzgodnieniu orientacji
$\lambda(e,f)=1$. Każda sfera ma
trywialną wiązkę normalną, więc
$\mu(e)=\mu(f)=0$. Jasne pasy pokazują
jedynie przekrój otoczeń, a nie całe sfery.}
\label{fig:para-hiperboliczna-22}
\end{figure}

\subsection*{Przesunięcie uchwytu jako zmiana bazy}
Stwierdzenie~\ref{prop:slide-18} opisało geometryczny
ruch dwóch uchwytów tego samego indeksu:
$e_i$ zostaje zastąpione przez $e_i\pm e_j$.
W wymiarze środkowym wpływ tego ruchu na
rachunek przecięć można wypisać bez kolejnego
twierdzenia geometrycznego. Jeżeli kolumny
macierzy $P$ wyrażają nową bazę w starej,
to macierz formy w nowej bazie wynosi
\begin{equation}\label{eq:slide-forma-22}
 A'=P^{T}AP,\qquad
 \mu(e_i\pm e_j)
   =\mu(e_i)+\mu(e_j)
      \pm[\lambda(e_i,e_j)].
\end{equation}
Pierwsza równość pochodzi z dwuliniowości
$\lambda$ nad $\Z$; druga z
\eqref{eq:quad-add-22} i
$\mu(-e_j)=\mu(e_j)$.
Na przykład dla symetrycznej pary hiperbolicznej
$A=\left(\begin{smallmatrix}0&1\\1&0\end{smallmatrix}\right)$
przesunięcie $e'=e+f$ daje
\[
 A'=\begin{pmatrix}2&1\\1&0\end{pmatrix},
 \qquad \mu(e')=1.
\]
Pojawienie się niezerowego wyrazu na przekątnej
nie tworzy nowej przeszkody: zmieniła się baza
tej samej formy. Przy niebanalnej $\pi_1$ w
\eqref{eq:slide-forma-22} trzeba zastąpić
transpozycję transpozycją sprzężoną oraz
śledzić etykiety dróg; sama macierz całkowita
nie zawiera wtedy całej informacji.

\begin{stwierdzenie}[Dlaczego hiperboliczna para nie przeszkadza]
\label{prop:hiperboliczna-22}
W parze $(e,f)$ mamy $\lambda(e,f)=1$,
$\lambda(e,e)=\lambda(f,f)=0$ oraz
$\mu(e)=\mu(f)=0$. Formy hiperboliczne
tworzą neutralne składniki przy stabilizacji.
Ich sygnatura, gdy $\epsilon=+1$ i $\Lambda=\Z$,
jest równa zeru.
\end{stwierdzenie}
\begin{proof}
Pierwsze równości odczytujemy z macierzy.
Na prostą generowaną przez $e$ forma
i funkcja kwadratowa znikają, a jej
ortogonalne dopełnienie wyznacza warunek
$\lambda(ae+bf,e)=\epsilon b=0$, więc
to znowu $\langle e\rangle$.
Dla $\epsilon=+1$ macierz ma wartości własne
$+1$ i $-1$, stąd jej sygnatura wynosi zero.
Suma wielu par ma lagranżjan będący sumą
odpowiednich prostych.
\end{proof}

\begin{figure}[htbp]
\centering
\begin{tikzpicture}[>=Latex,font=\small,line cap=round]
 \shade[inner color=white,outer color=manifoldblue!18,
        rounded corners=8pt] (-2.7,-1.08) rectangle (2.7,1.08);
 \draw[manifoldblue!60!black,rounded corners=8pt]
       (-2.7,-1.08) rectangle (2.7,1.08);
 \draw[manifoldblue!70!black,densely dashed]
       (-2.30,.72)--(2.30,.72);
 \draw[manifoldred!82!black,very thick,->]
       (-2.06,0)--(2.10,0)
       node[near end,below] {$e$};
 \draw[manifoldgreen!75!black,very thick,->]
       (0,-.81)--(0,.80)
       node[near end,right] {$f$};
 \fill[manifoldgold!85!black] (0,0) circle (2.9pt);
 \node[anchor=west] at (3.03,.48)
       {$\lambda(e,f)=1$};
 \node[anchor=west] at (3.03,-.03)
       {$\mu(e)=\mu(f)=0$};
 \node[anchor=west] at (3.03,-.55)
       {$\langle e\rangle=\langle e\rangle^\perp$};
\end{tikzpicture}
\caption{Lokalny schemat pary hiperbolicznej.
Sfery reprezentujące $e$ i $f$ spotykają się poprzecznie
w jednym punkcie, a każdą z nich osobno można
obramować bez samoprzecięć. Wybór kierunku $e$
wyznacza lagranżjan; prostokąt jest tylko fragmentem
otoczenia przecięcia, nie całą rozmaitością.}
\label{fig:hiperboliczna-22}
\end{figure}

\begin{przyklad}[Produkt sfer z poprzedniego rozdziału]
\label{prz:produkt-L-22}
W Przykładzie~\ref{prz:forma-produkt-21}
klasy $e=[S^k\times\{*\}]$ i
$f=[\{*\}\times S^k]$ mają macierz $H_{(-1)^k}$.
Ich wiązki normalne są trywialne, a reprezentanci
nie mają samoprzecięć, toteż $\mu(e)=\mu(f)=0$.
To przykład jądra \emph{niezerowego jako moduł},
które daje \emph{zerową przeszkodę}:
można usunąć parę za pomocą odpowiednich chirurgii.
\end{przyklad}

\section{Co oznacza grupa \texorpdfstring{$L$}{L}}
\label{sec:grupy-L-22}

W wymiarze $2k$ sumujemy formy kwadratowe
ortogonalnie. Dwie formy uważa się za równoważne
w sensie Witta, jeśli ich ortogonalna różnica,
po ewentualnym dodaniu form hiperbolicznych,
ma lagranżjan. W wersji $L^s$ moduły mają
preferowane bazy stabilne, a dla lagranżjanu
$L\subset P$ odwzorowanie $P/L\to L^*$
wyznaczone przez $\lambda$ ma być \emph{prostym}
izomorfizmem względem tych baz. Jest to algebraiczny warunek
zerowania różnicy klas. Klasy z dodawaniem
przez sumę ortogonalną tworzą \emph{kwadratową grupę
chirurgiczną} $L_{2k}(\Lambda)$, z odpowiednią
dekoracją opisaną poniżej. Wersję opartą na bazach
i prostych równoważnościach oznaczymy
$L^s_{2k}(\Lambda,w)$. Warunek „prosta” dotyczy
również dualności formy: izomorfizm do modułu dualnego
ma zerową torsję w $\operatorname{Wh}(\pi)$.
Wersja $L^h_{2k}$ używa zwykłej równoważności
homotopijnej. Dla $\pi=1$ te rozróżnienia znikają.

W nieparzystym wymiarze jądro nie jest opisywane
pojedynczą formą środkową. Model algebraiczny ma
postać \emph{formacji}: hiperboliczny moduł
kwadratowy wraz z dwoma lagranżjanami. Ich
przecięcie i iloraz przez ich sumę mierzą dwa
sąsiednie jądra mapy; wyprowadzimy to niżej.
Formacje dzielimy przez stabilizację i relację
wyznaczoną przez kobordyzm formacji; otrzymujemy
$L^s_{2k+1}(\Lambda,w)$ lub $L^h_{2k+1}(\Lambda,w)$.
W pełnej definicji obejmującej wszystkie wymiary
są to klasy kobordyzmu skończonych
\emph{kwadratowych kompleksów Poincarégo}
nad $\Lambda$. Konstrukcja formacji jest ich
nieparzystowym opisem po usunięciu homotopijnych
składników poniżej środka. Nie utożsamiamy
nieparzystej grupy z grupą Witta macierzy z poprzedniej sekcji.

\begin{uwaga}[Nie mylmy dwóch słów „torsja”]
W $\operatorname{Wh}(\pi)$ torsja oznacza torsję
\emph{równoważności łańcuchowej} z
Rozdziału~\ref{ch:s-kobordyzm}; w grupie $L$
klasyfikujemy formy i kwadratowe kompleksy.
To, że algebraiczne $K$- i $L$-teoria korzystają
z tych samych modułów nad $\Z[\pi]$, nie czyni
ich przeszkód tym samym obiektem.
\end{uwaga}

\section{Przypadek prostej spójności: sygnatura i Arf}
\label{sec:obliczenia-LZ-22}

Dla $\pi=1$ grupy kwadratowe w wysokowymiarowej
teorii chirurgii mają czterookresowy wzór
\begin{equation}\label{eq:L-Z-22}
 L_{4j}(\Z)\cong\Z,\qquad
 L_{4j+1}(\Z)=0,\qquad
 L_{4j+2}(\Z)\cong\Z/2,\qquad
 L_{4j+3}(\Z)=0.
\end{equation}
Wzór jest twierdzeniem klasyfikacyjnym teorii $L$;
poniżej wyjaśniamy jego dwa niezerowe inwarianty
i sprawdzamy konkretne formy. Nie wyprowadzamy
samych obliczeń grup nieparzystych z rachunku
dwóch macierzy: wymagałoby to teorii formacji.

\subsection*{Wymiary podzielne przez cztery}
Niech $n=4j$, $k=2j$. Środkowa forma jest
parzysta, symetryczna i unimodularna.
Jej \emph{sygnatura} to liczba dodatnich minus
liczba ujemnych wartości własnych nad $\R$.
Klasyczny rachunek form całkowitych daje
\[
 \sigma(f,b)=\frac{\operatorname{sign}\lambda_K}{8}
       \in L_{4j}(\Z)\cong\Z.
\]
Podzielność sygnatury przez $8$ dla parzystej
formy unimodularnej oraz to, że powyższy wzór
daje izomorfizm grup Witta, to nietrywialne
twierdzenia o formach kwadratowych. Nie wynikają
z samego zdiagonalizowania formy nad $\R$.
Sygnatura hiperbolicznej pary jest zerowa, jak
wykazaliśmy w Stwierdzeniu~\ref{prop:hiperboliczna-22}.

\begin{przyklad}[Forma $E_8$ ma niezerową klasę]
\label{prz:E8-L-22}
Niech $E_8$ będzie macierzą symetryczną $8\times8$:
na przekątnej ma $2$, przy krawędziach
\[
 (1,3),(3,4),(4,5),(5,6),(6,7),(7,8),(2,4)
\]
ma $-1$, a poza nimi zera. Schemat jej niezerowych
wyrazów pozadiagonalnych ma postać
\[
\begin{tikzpicture}[x=.72cm,y=.65cm,baseline=(current bounding box.center)]
 \draw[manifoldblue!68!black,thick]
   (0,0)--(1,0)--(2,0)--(3,0)--(4,0)--(5,0)--(6,0)
   (2,0)--(2,1);
 \foreach \num/\x/\y in {1/0/0,3/1/0,4/2/0,5/3/0,
                            6/4/0,7/5/0,8/6/0,2/2/1}
   \node[circle,draw=manifoldblue!75!black,fill=white,
         inner sep=1.5pt,font=\scriptsize] at (\x,\y) {$\num$};
\end{tikzpicture}
\]
Jej wiodące minory główne wynoszą kolejno
$2,4,6,5,4,3,2,1$. Można to sprawdzić, kolejno
rozwijając wyznacznik po liściu narysowanego
drzewa: jeśli usuwamy liść $v$ o jedynym
sąsiedzie $u$, wyznacznik jest
$2\det A_{G-v}-\det A_{G-\{u,v\}}$.
Zatem kryterium Sylvestera daje dodatnią
określoność, a ostatni minor $1$ daje
unimodularność. Dla całkowitego wektora $x$
liczba $x^TE_8x$ jest parzysta, bo składniki
pozadiagonalne występują podwojone. Definiujemy
$\mu(x)=x^TE_8x/2$. Otrzymujemy formę
kwadratową o sygnaturze $8$, a więc
$[E_8]=1\in L_{4j}(\Z)$.
Jest to \emph{przykład algebraiczny}; realizowanie
zadanej formy przez konkretną gładką mapę
normalną wymaga osobnego twierdzenia realizacyjnego.
\end{przyklad}

\subsection*{Wymiary $4j+2$}
Niech $k=2j+1$. Redukcja $P$ modulo $2$
daje przestrzeń symplektyczną
$V=P/2P$ nad $\mathbb F_2$ z formą
$b=\lambda\bmod2$. Funkcja $q=\mu\bmod2$
spełnia
\begin{equation}\label{eq:Arf-q-22}
 q(x+y)=q(x)+q(y)+b(x,y).
\end{equation}
Dla bazy symplektycznej $(e_1,f_1,\dots,e_r,f_r)$
definiujemy
\[
 \operatorname{Arf}(q)
       =\sum_{i=1}^r q(e_i)q(f_i)\pmod2.
\]

\begin{stwierdzenie}[Niezależność od bazy]
\label{prop:Arf-baza-22}
Podany wzór na $\operatorname{Arf}(q)$ nie zależy
od wyboru bazy symplektycznej. Suma ortogonalna
dodaje inwarianty Arfa.
\end{stwierdzenie}
\begin{proof}
Policzmy liczbę niezależną od bazy:
\[
 S(q)=\sum_{x\in V}(-1)^{q(x)}.
\]
Z~\eqref{eq:Arf-q-22} i rozkładu na ortogonalne
płaszczyzny $\langle e_i,f_i\rangle$ wynika,
że $q$ na jednej płaszczyźnie przyjmuje na
$ae_i+bf_i$ wartość
$a q(e_i)+bq(f_i)+ab$, gdzie $a,b\in\mathbb F_2$.
Sumując cztery składniki, dostajemy
\[
 \sum_{a,b\in\mathbb F_2}
  (-1)^{a q(e_i)+bq(f_i)+ab}
     =2(-1)^{q(e_i)q(f_i)}.
\]
Iloczyn tych równości daje
$S(q)=2^r(-1)^{\sum_iq(e_i)q(f_i)}$.
Lewa strona nie używa bazy, więc prawa
wyznacza jednoznacznie Arfa. Dla sumy
ortogonalnej sumy $S$ mnożą się, dlatego
wykładniki znaków, czyli inwarianty Arfa,
dodają się modulo $2$.
\end{proof}

\begin{przyklad}[Dwie formy z tą samą macierzą przecięć]
\label{prz:Arf-dwie-22}
Niech $P=\Z e\oplus\Z f$ i
$\lambda(e,f)=1$, $\lambda(f,e)=-1$,
$\lambda(e,e)=\lambda(f,f)=0$.
Gdy $\mu(e)=\mu(f)=0$, para jest hiperboliczna
i $\operatorname{Arf}(q)=0$. Gdy
$\mu(e)=\mu(f)=1\in\Z/2$, otrzymujemy
$\operatorname{Arf}(q)=1$.
W drugim przypadku wszystkie trzy niezerowe
wektory płaszczyzny nad $\mathbb F_2$ mają
wartość $q=1$, bo
$q(e+f)=1+1+b(e,f)=1$.
Nie ma więc nawet jednowymiarowego
lagranżjanu z $q=0$. Macierz $\lambda$
jest w obu przykładach \emph{identyczna};
o przeszkodzie decyduje dodatkowa informacja
samoprzecięciowa.
\end{przyklad}

W tym wymiarze izomorfizm
$L_{4j+2}(\Z)\cong\Z/2$ przydziela formie
klasę $\operatorname{Arf}(q)$. Dowód, że
\emph{każda} forma z Arfem zerowym jest
stabilnie hiperboliczna, jest algebraiczną
częścią wspomnianego twierdzenia klasyfikacyjnego.

\section{Przeszkoda normalnej mapy i twierdzenie Walla}
\label{sec:tw-Walla-22}

Rozważmy odwzorowanie normalne stopnia jeden
$(f,b):M^n\to X^n$ z
Sekcji~\ref{sec:jadra-chirurgii-21}.
Zakładamy, że $X$ jest zamkniętą gładką
rozmaitością lub skończonym kompleksem
Poincarégo z wybranymi danymi normalnymi.
Po chirurgiach \emph{poniżej} środka można
w dogodnych warunkach doprowadzić $f$ do
wysokiej spójności. W wymiarze parzystym
$n=2k$ pozostałe jądro nad $\Z[\pi]$
z formą $(\lambda_K,\mu_K)$ reprezentuje
\[
 \theta(f,b)=[K,\lambda_K,\mu_K]
       \in L^s_{2k}(\Z[\pi],w).
\]
Bez wcześniejszych chirurgii tę samą klasę
można zdefiniować na całym kompleksie łańcuchów
jako klasę kwadratowego kompleksu Poincarégo.
To ostatnie ujęcie obejmuje też wymiar nieparzysty.
Klasa nie zależy od wybranej kolejności
chirurgii: ślady tych zmian dają kobordyzm
odpowiednich kompleksów kwadratowych.

\begin{twierdzenie}[Twierdzenie Walla o przeszkodzie chirurgicznej]
\label{thm:Wall-L-22}
\index{twierdzenie Walla@Twierdzenie Walla}
Niech $M^n$ i $X^n$ będą spójnymi, zamkniętymi,
gładkimi rozmaitościami, $n\geq6$, i niech
$(f,b):M\to X$ będzie odwzorowaniem normalnym
stopnia jeden z Definicji~\ref{def:normal-map-21},
z danymi normalnymi nad pewną stabilną wiązką
$\xi\to X$ jak w~\eqref{eq:normal-map-21}.
Oznaczmy $\pi=\pi_1(X)$ i przez $w$ charakter
orientacji. Wtedy przeszkoda
$\theta(f,b)\in L^s_n(\Z[\pi],w)$ jest niezmiennikiem
normalnej klasy kobordyzmu.
Równość $\theta(f,b)=0$ zachodzi wtedy i tylko wtedy,
gdy $(f,b)$ jest normalnie kobordantne z prostą
równoważnością homotopijną $f':M'\to X$.
W wersji zwartej z brzegiem zakładamy ponadto,
że odwzorowanie na brzegu jest prostą równoważnością
homotopijną i wykonujemy kobordyzm oraz chirurgie
relatywnie do brzegu. Wersja dla zwykłych
równoważności używa dekoracji $h$.
\end{twierdzenie}

\noindent\textbf{Zakres wyprowadzenia.}
Rozpiszemy argument dostateczności dla $\pi=1$
i $n=2k\geq6$, a część rachunku przecięć
i lagranżjanu także nad $\Z[\pi]$.
Pierwszy model formacji wyjaśni przypadek
nieparzysty. Pozostałe wejścia geometryczne
pochodzą z
\href{https://webhomes.maths.ed.ac.uk/~v1ranick/books/scm.pdf}
{C.~T.~C. Walla, \emph{Surgery on Compact Manifolds},
§§1, 5--6, szczególnie twierdzenia 1.2, 5.2,
5.6, 6.4 oraz lemat 5.7}.
Lemat 5.3 o prostym lagranżjanie dowiedziemy
poniżej. Pełna wersja dla dowolnej $\pi$
i nieparzystego $n$ nadal wymaga dalszych
lematów geometrycznych i relacji formacji.

\noindent\textbf{Co dokładnie oddziela poniższy rachunek
od pełnego twierdzenia?}
Są to cztery odrębne kroki, a nie samo
wydłużenie rachunku form nad $\Z$:
\begin{enumerate}[label=(\arabic*)]
\item \emph{Chirurgie poniżej środka.}
  Poniżej konstruujemy zgodne obramowanie
  dla pojedynczej klasy w zakresie $2r<n$.
  Skończona indukcja po względnych komórkach
  i kontrola nowych jąder jest treścią
  twierdzenia 1.2 Walla.
\item \emph{Realizacja w środku.}
  Wyliczymy etykiety przecięć i zbudujemy
  bazę hiperboliczną. Osadzenie sfery w środku
  z warunku $\mu=0$ wymaga obramowanych dysków
  Whitneya; dla wielu sfer trzeba usuwać
  również ich wzajemne przecięcia zgodnie
  z formą $\lambda$ i udoskonaleniem $\mu$;
  używamy tu twierdzenia 5.2 Walla.
\item \emph{Prosta równoważność i niezmienniczość.}
  Zmiana kolejności chirurgii musi dawać
  tę samą klasę przeszkody. Dla ogólnej
  grupy podstawowej wymaga to bazowanych
  kompleksów nad $\Z[\pi]$, kontroli torsji
  Whiteheada i analizy kobordyzmu normalnego.
  Poniżej wypiszemy ciąg dokładny
  kobordyzmu, lecz jego przygotowanie
  wymaga twierdzenia 1.4 i lematu 5.7 Walla.
\item \emph{Wymiar nieparzysty.}
  Jądra dwóch sąsiednich wymiarów są mierzone
  przecięciem i ilorazem dwóch lagranżjanów
  formacji. Poniżej wyprowadzimy jej ciąg
  dokładny; relacja stabilna i realizacja
  odpowiadających jej ruchów uchwytów
  pozostają wejściem z §6 Walla.
\end{enumerate}
Dlatego nie nazywamy poniższego wyprowadzenia
pełnym dowodem twierdzenia w dowolnej grupie
podstawowej i każdym wymiarze.

\begin{lemat}[Algebraiczny efekt usunięcia jednej pary]
\label{lem:para-Wall-22}
Niech $(K,\lambda,\mu)$ będzie sumą $r$ płaszczyzn
hiperbolicznych nad $\Z$, z bazą
$e_1,f_1,\ldots,e_r,f_r$ taką, że
$\lambda(e_i,f_j)=\delta_{ij}$,
$\lambda(e_i,e_j)=\lambda(f_i,f_j)=0$
oraz $\mu(e_i)=\mu(f_i)=0$.
Wówczas
\[
 e_r^\perp/\Z e_r
 \cong\langle e_1,f_1,\ldots,e_{r-1},f_{r-1}\rangle
\]
jako formy kwadratowe.
\end{lemat}
\begin{proof}
Zapiszmy
$x=\sum_i(a_i e_i+b_i f_i)$.
Z definicji bazy hiperbolicznej
$\lambda(e_r,x)=b_r$. Stąd $x\in e_r^\perp$
dokładnie wtedy, gdy $b_r=0$.
Pozostały współczynnik $a_r$ znika po
przejściu do ilorazu przez $\Z e_r$.
Pary o indeksach $i<r$ są ortogonalne
do $e_r,f_r$, więc ich $\lambda$ i $\mu$
przechodzą do ilorazu bez zmiany.
\end{proof}

\begin{lemat}[Od prostego lagranżjanu do bazy hiperbolicznej]
\label{lem:lagr-hip-Z-22}
Niech $(P,\lambda,\mu)$ będzie nieosobliwą
formą $\epsilon$-kwadratową nad
$\Lambda=\Z[\pi]$, $\epsilon=(-1)^k$.
Załóżmy, że $P$ jest wolny z preferowaną bazą,
a $L\subset P$ jest wolnym lagranżjanem z bazą
$(e_1,\ldots,e_r)$, która przedłuża się do
preferowanej bazy $P$. Niech ponadto
$P/L\to L^*$, $[x]\mapsto\lambda(-,x)|_L$,
będzie prostym izomorfizmem.
Wtedy istnieje preferowana baza
$(e_1,f_1,\ldots,e_r,f_r)$, w której
$\lambda(e_i,f_j)=\delta_{ij}$ oraz
$\lambda(e_i,e_j)=\lambda(f_i,f_j)=0$ i
$\mu(e_i)=\mu(f_i)=0$.
\end{lemat}
\begin{proof}
Wybierzmy podniesienia $f_j^0\in P$ dualnej
bazy $L^*$, zatem
$\lambda(e_i,f_j^0)=\delta_{ij}$.
Jądrem $P\to L^*$ jest $L^\perp=L$,
więc klasy $f_j^0$ są bazą $P/L$.
Założenie o prostocie zapewnia, że
$(e_i,f_j^0)$ reprezentuje preferowaną
klasę baz $P$.

Połóżmy $B_{ij}=\lambda(f_i^0,f_j^0)$
i wybierzmy $m_i\in\Lambda$ reprezentujące
$\mu(f_i^0)$. Zdefiniujmy
\begin{equation}\label{eq:poprawka-lagr-Wall-22}
 f_i=f_i^0-m_i e_i-
       \sum_{j>i}B_{ij}e_j.
\end{equation}
Dodajemy wyłącznie wektory z $L$, toteż
$\lambda(e_i,f_j)=\delta_{ij}$.
Dla $i<j$ składnik $-B_{ij}e_j$ w $f_i$
znosi pierwotne parowanie:
\[
 \lambda(f_i,f_j)=B_{ij}-B_{ij}=0.
\]
Przy $i>j$ używamy hermitowskości.
Ponieważ $\lambda(f_i^0,e_j)=0$ dla $j>i$
i $\lambda(L,L)=0$, równanie
\eqref{eq:quad-add-22} daje
\[
 \mu(f_i)=[m_i]+[\lambda(f_i^0,-m_i e_i)]
   =[m_i-\epsilon\bar m_i]=0.
\]
Równanie~\eqref{eq:quad-diag-22} daje teraz
$\lambda(f_i,f_i)=0$.
Zmiana $f_i^0\mapsto f_i$ jest elementarną
zmianą bazy i nie zmienia jej prostej klasy.
To dowodzi tezy również przy nieprzemiennym
$\Lambda$; dla $\Lambda=\Z$ odzyskujemy
poprzedni rachunek parzystych macierzy
i poprawki Arfa.
\end{proof}

\begin{stwierdzenie}[Przeciwna forma daje odwrotność]
\label{prop:odwrotna-forma-Wall-22}
Jeśli $(P,\lambda,\mu)$ jest prostą
nieosobliwą formą $\epsilon$-kwadratową
na wolnym bazowanym module, to
\[
 (P,\lambda,\mu)\oplus(P,-\lambda,-\mu)
\]
jest formą hiperboliczną. Stąd jej druga
składowa reprezentuje klasę przeciwną
w $L^s_{2k}(\Lambda,w)$.
\end{stwierdzenie}
\begin{proof}
Weźmy diagonalny podmoduł
$\Delta=\{(x,x):x\in P\}\subset P\oplus P$.
Na $\Delta$ obie formy znoszą się:
$\lambda((x,x),(y,y))=0$ oraz
$\mu((x,x))=0$. Jest to bezpośredni
składnik z bazą $(e_i,e_i)$, jeśli
$(e_i)$ jest bazą $P$. Iloraz
$(P\oplus P)/\Delta$ identyfikujemy
z $P$ przez $[(x,y)]\mapsto x-y$.
Parowanie z $\Delta$ jest, z dokładnością
do znaku, izomorfizmem dualności
$P\to P^*$ wyznaczonym przez $\lambda$.
Jest proste z założenia. Lemat
\ref{lem:lagr-hip-Z-22} pokazuje więc,
że $\Delta$ daje bazę hiperboliczną.
\end{proof}

\begin{lemat}[Dlaczego obramowanie istnieje poniżej środka]
\label{lem:ramy-ponizej-srodka-Wall-22}
Niech $(f,b):M^{2k}\to X^{2k}$ będzie
odwzorowaniem normalnym, $k\geq3$, a
$\alpha\in\pi_{r+1}(f)$, gdzie $1\leq r\leq k-1$.
Klasa $\alpha$ ma reprezentanta będącego
osadzoną $r$-sferą z obramowaniem zgodnym
z danymi normalnymi. Chirurgia na tej sferze
zabija $\alpha$ i nie przywraca jąder
w stopniach mniejszych od $r+1$.
\end{lemat}
\begin{proof}
Reprezentant klasy względnej jest parą
$u:S^r\to M$ i $F:D^{r+1}\to X$ z
$F|_{S^r}=f\circ u$. Ponieważ
$2r<2k$, ogólne położenie pozwala
zastąpić $u$ homotopicznym osadzeniem:
oczekiwany wymiar jego samoprzecięć
wynosi $2r-2k<0$.

Połóżmy $q=2k-r$. Stabilne dane normalne
$b$ utożsamiają $u^*\nu_M$ z
$(f\circ u)^*\xi$ po stabilizacji.
Ta ostatnia wiązka przedłuża się przez
$F$ nad dyskiem, więc jest stabilnie
trywialna. Ponieważ
$u^*TM\oplus u^*\nu_M$ jest stabilnie
trywialna oraz $TS^r\oplus\R$
jest trywialna, wzór
$u^*TM=TS^r\oplus\nu_u$ pokazuje,
że również $\nu_u$ ma zgodne stabilne
obramowanie. Ponieważ
$q\geq r+2$, można usunąć stabilizację.
Istotnie, $O(q+1)/O(q)\cong S^q$;
długi ciąg homotopii i $(q-1)$-spójność
$S^q$ pokazują, że $O(q)\to O$
jest izomorfizmem na $\pi_j$ dla
$j\leq r$. Dotyczy to zarówno przeszkody
do trywialności $\nu_u$ nad $S^r$,
jak i różnicy dwóch jej obramowań.
Wybieramy rzeczywiste obramowanie.
Jeżeli po zestawieniu go z różniczką
$du:TS^r\to u^*TM$ jego stabilna klasa
różni się od ramy wyznaczonej przez $b$
i $F$ o element $\pi_r(O)$, podnosimy
ten element do $\pi_r(O(q))$ i skręcamy
obramowanie o jego odwrotność.
Otrzymujemy dokładnie zadaną stabilną
klasę ramy, a nie tylko dowolną
trywializację $\nu_u$.

Dysk $F$ i zgodne obramowanie przedłużają
normalne dane przez ślad, zgodnie z
Twierdzeniem~\ref{thm:przedluzenie-mapy-21}.
Ślad ma względem $M$ jeden uchwyt indeksu
$r+1$, który zabija $\alpha$.
Względem nowego brzegu indeks wynosi
$2k-r\geq k+1>r+1$; dlatego ciąg
homotopii pary nie wprowadza nowych
jąder poniżej stopnia $r+1$.
\end{proof}

Przypadek $r=0$ odpowiada dołączeniu uchwytu
indeksu jeden wzdłuż dwóch punktów. Jeżeli
$M$ ma kilka składowych, takimi uchwytami
możemy najpierw je połączyć, przedłużając
mapę wzdłuż wybranych dróg w spójnym $X$.
Ramy obu dysków zaczepienia dobieramy
zgodnie z danymi normalnymi $b$ dokładnie
tak jak w przypadku $h_i=1$ poniżej.

\begin{stwierdzenie}[Skończona redukcja poniżej środka;
 szczególny przypadek twierdzenia 1.2 Walla]
\label{prop:skonczona-redukcja-Wall-22}
Niech $(f,b):M^{2k}\to X^{2k}$ będzie normalnym
odwzorowaniem zwartej rozmaitości gładkiej
bez brzegu,
$k\geq3$, a $X$ niech będzie spójną
przestrzenią o typie skończonego kompleksu
CW z danymi normalnymi wymaganymi dla $(f,b)$.
Istnieje skończony
ciąg chirurgii na sferach wymiarów mniejszych
od $k$, po którym otrzymujemy normalne
odwzorowanie $f':M'\to X$ będące
$k$-spójne. Wszystkie użyte uchwyty mają
indeks co najwyżej $k$.
\end{stwierdzenie}
\begin{proof}
Po ewentualnym połączeniu składowych zakładamy,
że $M$ jest spójna. Z
Twierdzenia~\ref{thm:cw-morse-17} wynika,
że $M$ ma skończony model CW: wybierzmy
równoważność $u:K\to M$ ze skończonym $K$.
Wybieramy również skończony model $L\simeq X$.
Przenosimy $f\circ u$ na mapę $K\to L$,
przybliżamy ją komórkowo i tworzymy
podwójny walec odwzorowań
$M\leftarrow K\to L$. Ma on skończoną
względną strukturę CW nad $M$ i typ
homotopijny zwykłego walca $f$;
oznaczmy ten model przez $C_f$. Wybieramy
jego strukturę komórkową tak, by
$M\hookrightarrow C_f$ było inkluzją,
$C_f\simeq X$, a jego względna struktura
komórkowa nie miała komórek wymiaru zero.
Względne wierzchołki łączymy z $M$
lasem krawędzi, którego każda składowa
spotyka $M$ w jednym wierzchołku,
i zwijamy ten las względem $M$. Wybieramy
równoważność $C_f\to X$, której złożenie
z inkluzją $M$ jest homotopijne z $f$;
dane normalne przenosimy wzdłuż tej homotopii.
Porządkujemy względne komórki wymiarów
co najwyżej $k$ według wymiaru i piszemy
\[
 C_0=M,\qquad C_i=C_{i-1}\cup e^{h_i},
 \qquad 1\leq h_i\leq k.
\]

Indukcyjnie budujemy ślad $W_i$ chirurgii
od $M$ do $M_i$ i równoważność homotopijną
$W_i\simeq C_i$ zgodną z odwzorowaniem do
$C_f$. Dla $i=0$ bierzemy $W_0=M\times I$.
Załóżmy, że zbudowaliśmy $W_{i-1}$.
Komórka $e^{h_i}$ ma odwzorowanie dołączające
$S^{h_i-1}\to C_{i-1}$ i klasę
charakterystyczną w
$\pi_{h_i}(C_i,C_{i-1})$. Dotychczasowe
uchwyty mają indeks co najwyżej $k$.
Po odwróceniu śladu względem $M_{i-1}$
ich indeksy wynoszą co najmniej
\[
 (2k+1)-k=k+1>h_i.
\]
Stąd para $(W_{i-1},M_{i-1})$ jest
$k$-spójna. Po zastąpieniu rozważanych
odwzorowań walcami odwzorowań i użyciu
$W_{i-1}\simeq C_{i-1}$ otrzymujemy,
ponieważ $h_i\leq k$, izomorfizm względnych
klas homotopii
\[
 \pi_{h_i}(C_f,M_{i-1})
 \xrightarrow{\ \cong\ }
 \pi_{h_i}(C_f,C_{i-1}).
\]
Podnosimy nim klasę dysku charakterystycznego
$e^{h_i}$, a nie tylko klasę jego brzegu.
Reprezentant podniesionej klasy składa się
ze sfery $u_i:S^{h_i-1}\to M_{i-1}$
i dysku $D^{h_i}\to C_f$ ją wypełniającego.
Po złożeniu z $C_f\simeq X$ otrzymujemy
klasę względną odwzorowania $f_{i-1}$.
Zgodność podniesienia z klasą
charakterystyczną mówi ponadto, że obraz
$u_i$ w $C_{i-1}$ jest homotopijny
z odwzorowaniem dołączającym komórkę.

Dla $h_i\geq2$
Lemat~\ref{lem:ramy-ponizej-srodka-Wall-22}
realizuje tę klasę przez osadzoną
obramowaną sferę $S^{h_i-1}$ i chirurgię
z przedłużeniem danych normalnych.
Dla $h_i=1$ sfera $S^0$ składa się
z dwóch punktów, a podniesiony dysk
jest drogą łączącą ich obrazy w $C_f$.
Wiązka normalna nad tą drogą jest
trywialna. Dobieramy małe dyski wokół
obu punktów z ramami zgodnymi ze stabilnymi
danymi $b$; jest to możliwe, bo
$\pi_0 O(2k)\to\pi_0 O$ jest izomorfizmem.
Dołączamy wzdłuż nich uchwyt indeksu jeden
i przedłużamy dane normalne, jak w przypadku
$r=0$ twierdzenia 1.1 Walla.

Nowy ślad powstaje z $W_{i-1}$ przez
dołączenie uchwytu indeksu $h_i$.
Zgodność klasy dołączającej z klasą
komórki i homotopijna niezmienniczość
dołączania komórek dają
$W_i\simeq C_i$ nad $C_f$.

Po ostatniej takiej komórce otrzymujemy
$C_j$. Para $(C_f,C_j)$ nie ma komórek
względnych wymiaru co najwyżej $k$,
więc jest $k$-spójna. Para $(W_j,M_j)$
też jest $k$-spójna, bo wszystkie jej
odwrócone uchwyty mają indeks co najmniej
$k+1$. Zatem złożenie
$M_j\to W_j\simeq C_j\hookrightarrow C_f$
jest $k$-spójne: pierwszy i ostatni człon
mają tę własność przez odpowiednie
ciągi homotopii par, a środkowy jest
równoważnością. Równoważność $C_f\simeq X$
pokazuje, że $f_j:M_j\to X$ jest
$k$-spójne.
Skończoność ciągu wynika ze skończoności
wybranego modelu CW.
\end{proof}

Dowód redukcji używa pojedynczego kroku
geometrycznego z Lematu~\ref{lem:ramy-ponizej-srodka-Wall-22}
i schematu indukcji z
\href{https://webhomes.maths.ed.ac.uk/~v1ranick/books/scm.pdf}
{Walla, §1, twierdzenie 1.2}. Ogólna wersja
obejmuje także inne wymiary i kategorię PL.

\begin{proof}[Rozwinięcie dowodu twierdzenia Walla w przypadku
 $\pi=1$, $n=2k\geq6$]
\textbf{Krok 1: chirurgie poniżej środka.}
Normalne dane $b$ pozwalają przedłużać
odwzorowanie nad śladami chirurgii, jak w
Twierdzeniu~\ref{thm:przedluzenie-mapy-21}.
Lemat~\ref{lem:ramy-ponizej-srodka-Wall-22}
konstruuje zgodne osadzenie i obramowanie
dla jednej klasy, a
Stwierdzenie~\ref{prop:skonczona-redukcja-Wall-22}
uzasadnia skończone powtarzanie tego kroku.
Po tych operacjach jedynym możliwym
jądrem chirurgii jest środkowe $K_k$.
Identyfikacja $K_k(M)\cong\pi_{k+1}(f)$
i prosta dualność Poincarégo dają
na nim nieosobliwą formę
$(K_k,\lambda,\mu)$ z
Sekcji~\ref{sec:kwadratowa-22}.
Wersja nad $\Lambda$ wymaga dodatkowo
preferowanej stabilnej bazy jądra:
to dokładnie lemat 5.1 Walla.
Podkreślamy, że warunek kwadratowy dotyczy
klas z \emph{normalnym} obramowaniem;
dla dowolnej sfery bez tych danych wzór
na samoprzecięcie ma dodatkowy składnik
Eulera wiązki normalnej (Wall, §5.2).

\textbf{Krok 2: przejście od klasy zerowej do
bazy hiperbolicznej.}
Równość $[K]=0$ w $L_{2k}(\Z)$
oznacza, że dla pewnych $r,s$
forma $K\oplus H_\epsilon^r$ jest
izomorficzna z $H_\epsilon^s$.
Równoważnie, po stabilizacji ma
lagranżjan; przejście od niego do
rzeczywistej bazy z parami $e_i,f_i$
rozpisuje Lemat~\ref{lem:lagr-hip-Z-22}.
Każdą dodawaną
płaszczyznę realizujemy \emph{lokalną}
chirurgią na trywialnej obramowanej
$(k-1)$-sferze: nowa rozmaitość jest
$M\#(S^k\times S^k)$, a dwie sfery
czynnikowe przecinają się raz z
odpowiednim znakiem i mają $\mu=0$.
Ślad dostarcza normalnego kobordyzmu.
To jest lemat 5.5 Walla; wyjaśnia, dlaczego
stabilizacja algebraiczna jest dozwolonym
ruchem geometrycznym. Otrzymujemy bazę
$e_i,f_i$ jak w Lemacie~\ref{lem:para-Wall-22}.

\textbf{Krok 3: jedna rzeczywista chirurgia
w środku.}
Warunek $\mu(e_r)=0$ i lemat o reprezentacji
Walla (§5.2), używający triku Whitneya
przy $k\geq3$, dają osadzoną obramowaną
sferę $S^k\subset M$ reprezentującą $e_r$
oraz dysk w $X$ wypełniający jej obraz.
Wykonujemy chirurgię nad $X$.
Jej ślad $W$ powstaje z $M\times I$
przez dołączenie uchwytu indeksu $k+1$.
Oznaczmy nowy brzeg przez $M'$ i piszmy
$K_i(W)$ dla jądra mapy śladu do $X\times I$.
Oba względne kompleksy uchwytowe mają po
jednym generatorze:
\[
 H_{k+1}(W,M;\Z)\cong\Z,\qquad
 H_k(W,M';\Z)\cong\Z.
\]
W pierwszym ciągu dokładnym mapa brzegu
$\Z\to K_k(M)$ wysyła $1$ na $e_r$.
Jest injektywna, gdyż partner $f_r$ daje
$\lambda(e_r,f_r)=1$, więc $e_r$ jest
prymitywny. Otrzymujemy
\[
 0\longrightarrow\Z e_r\longrightarrow K_k(M)
 \longrightarrow K_k(W)\longrightarrow0.
\]
Od strony $M'$ uchwyt ma indeks $k$.
Złożenie $K_k(M)\to K_k(W)\to
H_k(W,M';\Z)\cong\Z$ mierzy przecięcie
z $e_r$. Dla lewych modułów nad
$\Lambda$ odpowiada mu \emph{liniowy}
homomorfizm
\[
 \ell_{e_r}:K_k(M)\longrightarrow\Lambda,
 \qquad \ell_{e_r}(z)=\epsilon\lambda(z,e_r),
 \quad \epsilon=(-1)^k.
\]
Istotnie, $\ell_{e_r}(az)=a\ell_{e_r}(z)$,
a hermitowskość pokazuje, że jego jądrem
jest $e_r^\perp$. Przy $\Lambda=\Z$ jest to
ten sam wzór co $\lambda(e_r,z)$.
Ponieważ $\lambda(e_r,e_r)=0$, mamy
$\ell_{e_r}(e_r)=0$; mapa przechodzi więc
przez $K_k(W)=K_k(M)/\Z e_r$.
Drugi ciąg dokładny daje
\[
 0\longrightarrow K_k(M')\longrightarrow K_k(W)
 \xrightarrow{\,\ell_{e_r}\,}\Z
 \longrightarrow K_{k-1}(M')\longrightarrow0.
\]
Ponieważ $f_r$ przechodzi na $1$,
$K_{k-1}(M')=0$, a
$K_k(M')\cong e_r^\perp/\Z e_r$.
Forma i udoskonalenie kwadratowe przechodzą
na ten iloraz, gdyż $\mu(e_r)=0$.
W ogólnym przypadku wszystkie wystąpienia
$\Z$ w tych ciągach zastępuje $\Lambda$.
Bez partnera o przecięciu $1$ mógłby
powstać niezerowy $K_{k-1}(M')$.
Lemat~\ref{lem:para-Wall-22} usuwa dokładnie
parę $e_r,f_r$. Powtarzamy ten krok $r$ razy.

\textbf{Krok 4: otrzymanie równoważności.}
Po usunięciu wszystkich par $K_k=0$,
a inne jądra zostały usunięte w Kroku 1
lub znikają przez dualność. Zatem nowa
mapa stopnia jeden jest izomorfizmem
homologii. Ponieważ obie strony są
jednospójne i mają typ kompleksów CW,
twierdzenie Whiteheada
(Twierdzenie~\ref{thm:whitehead})
daje równoważność homotopijną.
Dla $\pi=1$ grupa Whiteheada jest zerowa,
więc jest to równoważność prosta.

\textbf{Konieczność i niezależność wyborów.}
Przeszkoda prostej równoważności wynosi
zero, bo jej kompleks jądra jest ściągalny.
Niech teraz $N^{2k+1}$ będzie normalnym
kobordyzmem między przygotowanymi mapami
$M_-$ i $M_+$. Obie mapy brzegowe
przygotowujemy do $k$-spójności, a mapę
$N\to X\times I$ czynimy $k$-spójną
chirurgiami relatywnymi. Odejmowanie uchwytów
usuwa $K_k(N,\partial N)$. Po lokalnej
stabilizacji preferowana stabilna baza
$K_{k+1}(N,\partial N)$ ma rzeczywistego
reprezentanta. Są to hipotezy lematu~5.7
Walla; ich osiągnięcie korzysta również
z jego twierdzenia~1.4. Wówczas otrzymujemy
bazowany ciąg
\[
 0\longrightarrow K_{k+1}(N,\partial N)
 \longrightarrow K_k(\partial N)
 \longrightarrow K_k(N)\longrightarrow0.
\]
Pierwszy i trzeci człon są dualne przez
dualność Poincarégo. Obraz pierwszego
jest więc kandydatem na lagranżjan w
\[
 K_k(\partial N)
 \cong K_k(M_+)\oplus K_k(M_-)^{-},
\]
gdzie minus oznacza przeciwną formę.
Dlaczego znikają na nim także $\lambda$
i $\mu$? Względną klasę reprezentuje się
obramowaną zanurzoną $(k+1)$-sferą z
usuniętymi dyskami. Przecięcia dwóch takich
powierzchni tworzą łuki, których końce na
brzegu mają przeciwne znaki i te same etykiety
dróg. Tak samo końce łuków samoprzecięć
znoszą się w $Q_\epsilon(\Lambda)$.
Dokładna realizacja reprezentantów,
odejmowanie uchwytów i prostota powyższego
ciągu są treścią twierdzenia 1.4 i lematu 5.7
Walla. Po ich zastosowaniu
Lemat~\ref{lem:lagr-hip-Z-22} daje formę
hiperboliczną na $K_k(\partial N)$,
a więc $[K_k(M_+)]=[K_k(M_-)]$.
Tu pozostaje jawne wejście geometryczne;
sam zapis ciągu nie dowodzi jego hipotez.
\end{proof}

\begin{figure}[htbp]
\centering
\begin{tikzpicture}[>=Latex,font=\small,
 box/.style={draw=manifoldblue!75!black,rounded corners=3pt,
 fill=manifoldblue!9,align=center,minimum height=1.05cm,
 text width=2.4cm}]
 \node[box] (a) at (0,0) {normalna mapa\\$M\to X$};
 \node[box] (b) at (3.25,0)
   {chirurgie poniżej\\środka};
 \node[box] (c) at (6.50,0)
   {forma jądra\\$(K,\lambda,\mu)$};
 \node[box] (d) at (9.75,0)
   {lagranżjan i\\chirurgie środkowe};
 \draw[->,thick,manifoldblue!75!black]
   (a)--(b);
 \draw[->,thick,manifoldblue!75!black]
   (b)--(c);
 \draw[->,thick,manifoldblue!75!black]
   (c)--(d);
 \draw[->,thick,manifoldgold!85!black]
   (d.south) -- ++(0,-1.30) -| (a.south);
 \node[fill=white,font=\scriptsize,align=center]
   at (4.85,-1.30) {zanik jądra daje równoważność};
\end{tikzpicture}
\caption{Kolejne etapy twierdzenia Walla w parzystym wymiarze.
Przejście od algebraicznego lagranżjanu do osadzonych
sfer wymaga lematów geometrycznych Walla.}
\label{fig:mechanizm-Wall-22}
\end{figure}

\subsection*{Dlaczego wymiar nieparzysty wymaga formacji}
Niech $n=2k+1\geq7$ i niech mapa normalna
została przygotowana chirurgiami poniżej
środka. Po stabilizacji wybieramy rozłączne
obramowane reprezentanty generatorów
$K_k(M)$, z otoczeniem
$U=\coprod(S^k\times D^{k+1})$.
To wybranie, wraz z kontrolą prostych baz,
jest geometrycznym krokiem z §6 Walla.
Połóżmy $C=\overline{M\setminus\operatorname{int}U}$.
Moduł $H=K_k(\partial U)$ ma formę
hiperboliczną: na każdej składowej brzegu
$S^k\times S^k$ dwie sfery czynnikowe
przecinają się raz. W obrazie odwzorowań
brzegu dostajemy dwa lagranżjany
\[
 L_0=\operatorname{im}\bigl(K_{k+1}(U,\partial U)
                  \to H\bigr),\qquad
 L_1=\operatorname{im}\bigl(K_{k+1}(C,\partial U)
                  \to H\bigr).
\]
Formacja jest trójką $(H;L_0,L_1)$.
Jej użycie nie oznacza, że dwa sąsiednie
jądra \emph{są} lagranżjanami; powstają
z ich wzajemnego położenia.

Algebraiczny ciąg
\begin{equation}\label{eq:formacja-ciag-Wall-22}
 0\longrightarrow L_0\cap L_1
 \longrightarrow L_1
 \longrightarrow H/L_0
 \longrightarrow H/(L_0+L_1)
 \longrightarrow0
\end{equation}
jest dokładny: środkowa strzałka wysyła
$z$ na $[z]$, więc jej jądrem jest
$L_0\cap L_1$, a jej obrazem
$(L_0+L_1)/L_0$.
Diagram Mayera--Vietorisa dla $M=U\cup C$
identyfikuje skrajne człony z
$K_{k+1}(M)$ i $K_k(M)$, odpowiednio.
Jeśli $L_0\oplus L_1=H$, oba te jądra
znikają; po zaniku niższych jąder i
przy $\pi_1=1$ twierdzenie Whiteheada
daje równoważność prostą.

Na przykład dla $\epsilon=-1$ weźmy
$H=\Z e\oplus\Z f$ z
$\lambda(e,f)=1$ oraz
$\mu(ae+bf)=ab\pmod2$.
Podmoduły $L_0=\Z e$ i
$L_1=\Z(e+2f)$ są lagranżjanami:
$\mu(e+2f)=0$, a wektor $e+2f$
jest prymitywny. Strzałka
$L_1\to H/L_0\cong\Z f$ z
\eqref{eq:formacja-ciag-Wall-22} jest
mnożeniem przez $2$. Jej cokernel to
$\Z/2$, zatem sam fakt istnienia dwóch
lagranżjanów nie oznacza, że obecne
jądro już znikło.

Przesunięcie uchwytu ma czytelny rachunek
w $H_\epsilon\oplus H_\epsilon$:
$e_1\mapsto e_1+e_2$ oraz
$f_2\mapsto f_2-f_1$, przy pozostałych
wektorach stałych, zachowuje wszystkie
parowania i $\mu$. Pełna relacja formacji
obejmuje ponadto stabilizację i formacje
brzegowe. Dlatego zerowa \emph{klasa} w
$L^s_{2k+1}(\Lambda,w)$ nie wymaga, aby
wybrane już teraz $L_0$ i $L_1$ były
dopełniające. Realizacja dozwolonych
ruchów przez chirurgie oraz niezależność
wyboru $U$ to treść
\href{https://webhomes.maths.ed.ac.uk/~v1ranick/books/scm.pdf}
{Walla, §6, twierdzenia 6.4 i 6.6};
algebraiczną relację formacji wyjaśnia
\href{https://www.maths.ed.ac.uk/~v1ranick/papers/ranicki.pdf}
{A.\ Ranicki, \emph{An Introduction to Algebraic Surgery}, §9}.

\begin{uwaga}[Gdzie kończy się argument z homologii]
\label{uwaga:homologia-L-22}
Zerowa zwykła homologia jądra nie wystarcza:
nad $\Z[\pi]$ mogą pozostać nietrywialne
klasy dróg. Nie wystarcza również nieosobliwość
formy: potrzeba lagranżjanu zgodnego z $\mu$.
Wreszcie istnienie \emph{homotopijnej}
równoważności nie rozstrzyga automatycznie
kwestii torsji Whiteheada potrzebnej dla
równoważności prostej.
\end{uwaga}

\section{Trzy rachunki i kolejny krok}
\label{sec:przyklady-L-22}

\begin{enumerate}[label=(\arabic*)]
\item Dla identyczności $\id:X\to X$ jądro
kompleksu mapy jest ściągalne, więc
$\theta(\id)=0$.
\item Dla normalnej mapy z
Przykładu~\ref{prz:produkt-L-22} środkowe
jądro ma rangę $2$, lecz jest hiperboliczne.
Dlatego jego klasa $L$ znika.
\item Algebraiczne formy $E_8$ oraz para
antysymetryczna z $\mu(e)=\mu(f)=1$
pokazują odpowiednio niezerową przeszkodę
sygnaturową i niezerową przeszkodę Arfa.
Niezerowość wyklucza istnienie końcowej
równoważności w klasie danej mapy normalnej.
\end{enumerate}

Rozdział~\ref{ch:chirurgia-geometryczna-21}
dał \emph{operację}, a ten rozdział daje
\emph{przeszkodę}. Do klasyfikowania
rozmaitości pozostaje zestawić normalne
niezmienniki, klasy struktur i grupy $L$
w ciągu dokładnym chirurgii. W takim ciągu
mapa do grupy $L$ mówi, które dane normalne
można zrealizować przez równoważność
homotopijną, a działanie poprzedniej grupy
$L$ opisuje niejednoznaczność wyniku.

\par\smallskip
{\footnotesize\noindent\emph{Dalsza lektura:}
\href{https://webhomes.maths.ed.ac.uk/~v1ranick/books/scm.pdf}
{C.~T.~C. Wall, \emph{Surgery on Compact Manifolds},
rozdziały o formach i przeszkodach},
\href{https://webhomes.maths.ed.ac.uk/~v1ranick/papers/ats.pdf}
{A.\ Ranicki, \emph{The Algebraic Theory of Surgery},
§7--8},
\href{https://link.springer.com/book/10.1007/978-3-031-56334-8}
{W.\ Lück i T.\ Macko, \emph{Surgery Theory:
Foundations}, rozdziały o teorii $L$}.\par}

\clearpage
